\newpage
\begin{center}
    \Huge Problem Statement
\end{center}

First, let us observe that the rigor of McDowell's system of proof is merely
illusory. In terms of classical logic, it constitutes an \emph{indirect proof}, whereby ``the truth of a proposition is established by demonstrating the falsehood of the assumption to its contrary.''\footnote{\emph{A Brief Dictionary of Logic}, Aleksandr Ivin, Dmitry Gorsky, \& Aleksandr Nikiforov, p. 81.} This method is not universal in and of itself, and there are a number of fundamental restrictions on its use. However, there is a more pressing issue here. In this case, the falsehood of the contrary assumption is not proved deductively, but is inferred as an inductive generalization, suggesting that this is actually a case of incomplete induction (a so-called ``hasty generalization'').

Moving from the logical aspect of our problem to the substantive aspect: it should be emphasized that McDowell's method of sequentially examining a list of alternatives is sound if, and only if, their combined whole exhaustively spans the entire range of logical possibilities (as McDowell in fact insists). Thus, in order to refute the ``McDowell-Gladkov System,'' the following is necessary and sufficient: I must propose at least one other hypothesis (naturally, a materialist one) that would provide a more consistent explanation than the ones previously refuted for the events surrounding Christ's crucifixion and resurrection.\footnote{I would like to note that in Damiano Damiani's outstanding film \emph{The Inquiry}, the emperor Tiberius sends a ``special investigator'' out to Jerusalem who attempts to do more or less the same thing (without much success, however). (K.E.)}

Thus, our problem is identical to the one that is solved in any classical (``English'') detective story, which is essentially like assembling a jigsaw puzzle. From a fixed number of pieces of a given shape (established facts), we need to construct a figure (alternate version) in such a way that the fragments fit together without gaps and every single one of them is used, even the most ``inconvenient'' ones. Just to be safe, let me reiterate: the only thing our version needs to be is \emph{internally consistent}; the question of its \emph{truth} or \emph{falsehood} is one that lies in a different plane altogether, and is simply irrelevant in this context.

Stating the question in this way allows us to, among other things, avoid discussing the historicity of Christ, the validity of treating the canonical Gospels as historical documents, and other such issues. It would be ridiculous for an amateur to try entangling himself in the problems that all kinds of specialists---historians, archeologists, philologists, linguists---have written vast piles of literature on. Therefore, by no means should we look at the text through the lens of academic Biblical studies; the primary focus of our investigation will be on \emph{literary characters}, not their historical prototypes (assuming they existed in the first place). I myself might define it as a continuation of the same tradition that produced one of my favorite books, \emph{The Riddle of Prometheus} by Lajos Mesterházi.

For this reason, we will take the authorship and timeline of the writing of the Gospels to be in full accordance with church canon; the problem of the continuity of the four canonical texts with the Logia of Matthew, the Gospel of Thomas, and the proto-Gospel ``Q'' (from the German \emph{Quelle}, ``source'') will be of no concern to us. My outline of the events will mostly follow that of Frederic Farrar's \emph{The Life of Christ}, a classic account that consolidates all the New Testament texts into a single narrative and is written from a conventional Christian perspective.

There are a number of places, however, where the various Evangelists' accounts seem to diverge quite fundamentally; in these cases, we cannot avoid discussing the source of these discrepancies. When approaching the accounts of Saints Matthew, Peter, Paul and John\footnote{While this list may seem strange, there is a traditional belief that the Gospels of Mark and Luke were based on the testimony of Saint Peter and Saint Paul respectively. (Z.B.)} that form the basis of their respective Gospels, we can think of them as... say, the testimony of four lodgers at a snowed-in Yorkshire estate where, by the will of Agatha Christie, a mysterious crime has taken place. In cases where the interpretation of the texts, or the historical reality underlying them, is contentious or dubious, I think it would be fair to give preference to McDowell's opinion (if he has expressed one).

While we're busy setting out our starting conditions, let us introduce an important constraint. In the ``materialist'' hypotheses discussed by McDowell, Christ and his associates are presented as being, to put it bluntly, either morons or crooks. And even though I myself am indifferent to religion, and quite cold towards the official Churches, these are assumptions that I strongly dislike. At the end of the day, no matter what the Church Fathers might say, every person has their own Christ. And while, from the point of view of orthodox dogmatics, \emph{my} Christ might be utterly monstrous (given that he bears all the ``birthmarks'' of liberal theology), under no circumstances is he a \emph{liar}. In any case, both Christ and the majority of his companions would soon go on to give their lives for their beliefs, something which ought to command some basic respect for them in and of itself. Therefore, when sorting through hypotheses that might serve to explain the various ``knots'' of events, we will proceed from the following assumption: we may only assume conscious deception on the part of Christ and the Apostles after all other possible explanations have been exhausted (the ``presumption of honesty''). Looking forward, I would like to note that such cases ultimately never occur---with just one caveat.

The English writer C.S. Lewis, a Christian philosopher and commentator, is
primarily known in our country as an author of didactic children's literature. In
one of his stories, a girl accidentally discovers a portal in an antique wardrobe that
leads to a magical land. When her older brother and sister hear her stories of her
visits, they begin to fear for their younger sister's sanity and ask their host, a
professor, for advice:

\begin{smallquote}
    ``Logic!'' said the Professor half to himself. ''Why don't they teach logic at these schools? There are only three possibilities. Either your sister is telling lies, or she is mad, or she is telling the truth. You know she doesn't tell lies and it is obvious that she is not mad. For the moment then and unless any further evidence turns up, we must assume that she is telling the truth.''\footnote{\emph{The Lion, the Witch and the Wardrobe}, p. 52}
\end{smallquote}

It is odd that Lewis' professor completely overlooks at least one other possibility: the fact that an honest and sane person can become sincerely deluded. The causes of such delusions are quite diverse. They could be, for example, various natural phenomena, from optical effects in the atmosphere (``flying saucers'') to systematic hallucinations caused by oxygen deprivation in high-altitude conditions (the ``abominable snowman,'' the ``Black Mountaineer'').\footnote{The Black Mountaineer is a ghost reportedly seen by some Soviet mountain climbers; a figure dressed entirely in black, he supposedly died in a mountaineering accident and now roams the mountains for all eternity. (Z.B.)} On the other hand, however---and this is more relevant to our case---any person can be the victim of a deliberate hoax.

What is significant here is that the perpetrator of a serious hoax not only needs to ensure that the staging itself is convincing, but also that it occurs in the presence of eyewitnesses whose reputation is nothing short of impeccable (as it would be useless to rely on the testimony of a fool or a known liar). There's a reason all kinds of specialists in telepathy and psychokinesis always want famous scientists to take part in their experiments, while categorically refusing to work their ``miracles'' in the presence of professional magicians. It is for precisely these reasons that our ``presumption of honesty'' does not mean that any fact reported by the Evangelists is necessarily true.